\section{Mixing Time Estimates}
\label{sec:mixing-time}

\subsection{Definition}

The \emph{total variation mixing time} of a Markov chain with transition
kernel $P$ is:
\[
  \tmix(\epsilon) = \min\left\{ t \geq 0 :
  \max_{x \in \Omega} \| P^t(x, \cdot) - \pi \|_{\mathrm{TV}} \leq \epsilon \right\},
\]
where $\|\mu - \nu\|_{\mathrm{TV}} = \frac{1}{2}\sum_x |\mu(x) - \nu(x)|$
is the total variation distance.
For practical purposes, $\epsilon = 0.01$ is a reasonable target.

\subsection{From R-hat to Mixing Time}

$\rhat < 1.1$ is a necessary but not sufficient condition for convergence
in total variation.
A rule of thumb from \citet{gelman2013} is that $\rhat < 1.1$ corresponds
roughly to $\| P^t(x, \cdot) - \pi \|_{\mathrm{TV}} < 0.1$ for the
scalar function $f$.
Full total-variation mixing at $\epsilon = 0.01$ may require several times
more steps.

For planning purposes, we use the $\rhat < 1.1$ criterion as a proxy for
mixing time, recognising that it provides a lower bound on the steps needed
for rigorous total-variation convergence.

\subsection{Empirical Estimates}

From Table~\ref{tab:mixing-times}, mixing time estimates for the six states
under the $\rhat < 1.1$ proxy are:

\begin{table}[h]
\centering
\caption{Estimated mixing time $\hat{t}_{\rm mix}$ (steps to $\rhat < 1.1$
         for Polsby-Popper) and steps per second for a 2024-generation
         workstation.}
\label{tab:mixing-time-est}
\begin{tabular}{lccc}
\toprule
State & $\hat{t}_{\rm mix}$ & Steps/sec & Wall time \\
\midrule
NC & 2{,}000 & 4.2 & $\approx$8 min \\
WI & 1{,}500 & 7.1 & $\approx$4 min \\
GA & 2{,}500 & 4.2 & $\approx$10 min \\
PA & 3{,}000 & 3.5 & $\approx$14 min \\
TX &10{,}000 & 1.6 & $\approx$104 min \\
CA &20{,}000 & 1.1 & $\approx$303 min \\
\bottomrule
\end{tabular}
\end{table}

For Texas and California, a single certified ensemble run exceeds two hours.
Parallel chains on a multi-core machine are essential.

\subsection{Comparison to Statutory Threshold}

The \cs{} algorithm with $T = 600$ typically completes in under 60 seconds
for any state (B.16 timing data).
A \recom{} ensemble certified to $\rhat < 1.1$ requires 4 minutes (WI) to
5 hours (CA) to collect sufficient samples.

This comparison illuminates the purpose of each method:
\begin{itemize}
\item \cs{} is appropriate for plan \emph{generation}: it produces one
      optimal plan quickly and deterministically.
\item Ensemble sampling is appropriate for plan \emph{evaluation}: it
      characterises the space of valid plans and can locate outliers.
\end{itemize}

The two methods are complementary, not competing.
